% ==================================================
% Fuentes, idioma y codificación
% ==================================================
\usepackage[T1]{fontenc}        % Acentos correctos en PDF
\usepackage[utf8]{inputenc}     % Archivos fuente en UTF-8
\usepackage[spanish]{babel}     % Español (títulos, fechas, etc.)
\decimalpoint                   % Usar punto decimal

% Fuente moderna (texto + mejor rasterizado)
\usepackage{lmodern}
%\usepackage{stix}

\usepackage[table]{xcolor} % Para pintar las celdas y filas
\usepackage{graphicx}      % Para el \resizebox
\usepackage{tikz}          % Para dibujar los cuadritos de la leyenda




% ==================================================
% Matemáticas
% ==================================================
\usepackage{amsmath, amssymb, mathtools}
% Opcionales de mates avanzadas (descomenta solo si los usas)
% \usepackage{physics}
%\usepackage{tensor}
% \usepackage{cancel}
\usepackage{empheq}             % Para enmarcar/destacar ecuaciones

% ==================================================
% Página, espaciado y párrafos
% ==================================================
\usepackage{setspace}

% ==================================================
% Cabeceras y pies de página
% ==================================================
%\usepackage{fancyhdr}
%\setlength{\headheight}{14.5pt}
%\pagestyle{fancy}


% ==================================================
% Tablas, múltiples columnas, etc.
% ==================================================
\usepackage{booktabs}           % Tablas bonitas
% \usepackage{multirow}         % Solo si usas celdas de varias filas
\usepackage{multicol}           % Texto en varias columnas

% ==================================================
% Colores, figuras y gráficos
% ==================================================
\usepackage{xcolor}
\usepackage{graphicx}
\graphicspath{{Figuras/}{logo/}}
\usepackage{float}
\usepackage{subcaption}
\usepackage{tikz}
\usetikzlibrary{positioning, shapes, matrix, fit, backgrounds, arrows.meta}
\pgfdeclarelayer{background}
\pgfdeclarelayer{foreground}
\pgfsetlayers{background,main,foreground}

\usepackage{pgfplots}
\pgfplotsset{
    width=10cm,
    compat=1.18    % versión moderna recomendada
}

% ==================================================
% Títulos de secciones (estética)
% ==================================================
\usepackage{titlesec}
\definecolor{gray75}{gray}{0.75}
\newcommand{\hsp}{\hspace{20pt}}

% Como extarticle/article NO tienen \chapter, formateamos \section:
\titleformat{\section}[hang]
  {\Large\bfseries}
  {\thesection\hsp\textcolor{gray75}{|}\hsp}
  {0pt}
  {\Large\bfseries}

% ==================================================
%  Descomentar para puntos en la tabla de contenidos
% ==================================================

%\usepackage{tocloft}
%\renewcommand{\cftsecleader}{\cftdotfill{\cftdotsep}}
%\sidecaptionvpos{figure}{tl}

% ==================================================
% Cajas de color / bloques (tcolorbox)
% ==================================================
%\usepackage{tcolorbox}
%\tcbuselibrary{skins,breakable}

% ==================================================
% Enlaces
% ==================================================
\usepackage[hidelinks]{hyperref}

%%%%%%%%%%%%%%%%%%%%%%%%%%%%%%%%%%%%%%%%%
%%%%%      Underline Eqs        %%%%%%%%%
%%%%%%%%%%%%%%%%%%%%%%%%%%%%%%%%%%%%%%%%%

\newsavebox\MBox
\newcommand\Cline[2][red]{{\sbox\MBox{$#2$}%
  \rlap{\usebox\MBox}\color{#1}\rule[-1.2\dp\MBox]{\wd\MBox}{0.5pt}}}
%%%%%%%%%%%%%%%%%%%%%%%%%%%%%%%%%%%%%%%%%%
%%%        CONSTANTES                %%%%%
%%%%%%%%%%%%%%%%%%%%%%%%%%%%%%%%%%%%%%%%%%
\newcommand\mom{\sqrt{\frac{m\omega}{2\hbar}}}
\newcommand\momsq{\frac{m\omega}{2\hbar}}
\newcommand\dad{\sqrt{\frac{m\omega\hbar}{2}}ie^{-\abs{z}^2/2}}
\newcommand\gus{\sqrt{\frac{\hbar}{2m\omega}}}
\newcommand\gussq{\frac{\hbar}{2m\omega}}
\newcommand\adag{\hat{a}^{\dagger}}
\newcommand\ah{\hat{a}}

\newcommand{\ylm}[2]{Y_{#1}^{#2}}
\newcommand{\fun}[3]{\Psi_{#1 #2 #3}(r,\theta,\phi)=R_{#1 #2}(r)Y_{#2}^{#3}(\theta,\phi)}
\newcommand{\funup}[3]{\Psi_{#1 #2 #3,\,+\frac{1}{2}}(r,\theta,\phi)=R_{#1 #2}(r)Y_{#2}^{#3}(\theta,\phi)}
\newcommand{\fundown}[3]{\Psi_{#1 #2 #3,\,-\frac{1}{2}}(r,\theta,\phi)=R_{#1 #2}(r)Y_{#2}^{#3}(\theta,\phi)}
%%%%%%%%%%%%%%%%%%%%%%%%%%%%%%%%%%%%%%%%%%%
%%%%     PALETA DE COLORES         %%%%%%%%
%%%%%%%%%%%%%%%%%%%%%%%%%%%%%%%%%%%%%%%%%%%
\definecolor{Verde}{HTML}{31B189}
\definecolor{Azul}{HTML}{7ACE67}
\definecolor{Naranja}{HTML}{FFC872}
% Paleta compartida con la presentación y el póster.
\definecolor{marinoCIMAT}{HTML}{0A1A2F}
\definecolor{vinoCIMAT}{HTML}{7A2237}
\definecolor{petroleo}{HTML}{3B5366}
\definecolor{fondoSlide}{HTML}{F4F5F7}
\definecolor{textoBody}{HTML}{1A1A1A}
\definecolor{textoSec}{HTML}{4A4A4A}
\definecolor{csvEconomico}{HTML}{A8C8E0}
\definecolor{csvComorbilidad}{HTML}{E89D9D}
\definecolor{csvTrabajoSocial}{HTML}{A8D5BA}
\definecolor{umapEconomico}{HTML}{378ADD}
\definecolor{umapComorbilidad}{HTML}{D85A30}
\definecolor{umapTrabajoSocial}{HTML}{1D9E75}
\definecolor{easyPositive}{HTML}{A8D5BA}
\definecolor{hardPositive}{HTML}{FFC966}
\definecolor{easyNegative}{HTML}{D8D8D8}
\definecolor{hardNegative}{HTML}{E89D9D}
\definecolor{bloqueCode}{HTML}{1C1C1E}
\definecolor{tokenCyan}{HTML}{5BD9C8}
\definecolor{tokenVerde}{HTML}{7EDFA5}
\definecolor{tokenVioleta}{HTML}{B4A3E8}
\definecolor{semVerde}{HTML}{4CAF50}
\definecolor{semAmarillo}{HTML}{FFC107}
\definecolor{semRojo}{HTML}{E53935}
\definecolor{Crema}{HTML}{FFE3B3}
\definecolor{AzulUV}{HTML}{1C36A1}
\definecolor{VerdeUV}{HTML}{00AA33}
\definecolor{RCimat}{HTML}{73243D}
\definecolor{GCimat}{HTML}{ACAAAB}


%%%%%%%%%%%%%%%%%%%%%%%%%%%%%%%%%%%%%%%%%%%%%%
%%%        Configuración bloques grises    %%%
%%%%%%%%%%%%%%%%%%%%%%%%%%%%%%%%%%%%%%%%%%%%%%
\usetikzlibrary{positioning, shapes}
\pgfdeclarelayer{background}
\pgfdeclarelayer{foreground}
\pgfsetlayers{background,main,foreground}
\usepackage{tcolorbox}
\tcbuselibrary{skins,breakable}
\newenvironment{myblock}[1]{%
    \tcolorbox[beamer,
    noparskip,breakable,
    colback=white,       % Fondo blanco
    colframe=black,      % Borde negro
    coltitle=black,      % Color del título en negro
    coltext=black,       % Texto negro
    sharp corners,       % Esquinas sin redondear
    boxrule=1pt,         % Grosor de las líneas del borde
    enhanced jigsaw,     % Mejora la estructura del cuadro
    before skip=10pt,    % Espacio antes del cuadro
    after skip=10pt]}    % Espacio después del cuadro
    {\endtcolorbox}


\newtcbox{\othermathbox}[1][]{nobeforeafter, tcbox raise base, enhanced, sharp corners, colback=white!10, colframe=AzulUV!30!black,drop fuzzy shadow, left=1em, top=1em, right=1em, bottom=1em}

% Syntax: \colorboxed[<color model>]{<color specification>}{<math formula>}
\newcommand*{\colorboxed}{}
\def\colorboxed#1#{%
  \colorboxedAux{#1}%
}
\newcommand*{\colorboxedAux}[3]{%
  % #1: optional argument for color model
  % #2: color specification
  % #3: formula
  \begingroup
    \colorlet{cb@saved}{.}%
    \color#1{#2}%
    \boxed{%
      \color{cb@saved}%
      #3%
    }%
  \endgroup
}
% ==================================================
\usepackage{csquotes}
\usepackage[backend=biber, style=numeric, sorting = none]{biblatex} % Estilo APA automático
\addbibresource{Bibliografia/referencias.bib} % Llama a tu archivo nuevo
